\documentclass[10pt,a4paper]{amsart}
\usepackage[margin=19mm]{geometry}
\usepackage{amsmath,amssymb,mathtools}
\usepackage[scr=rsfs,cal=euler]{mathalfa}
\usepackage{xfrac}
\usepackage[unicode,hidelinks,bookmarks=false]{hyperref}
\newcommand{\ie}{i.e.\ }
\newcommand{\cf}{cf.\ }
\newcommand{\resp}{respectively\ }
\newcommand{\Subsection}{\subsection}
\providecommand{\nobreakdash}{\nobreak\hbox{-}}
\setlength{\emergencystretch}{2em}
\multlinegap0pt
\usepackage{bm}
\usepackage{bbm}
\usepackage{dsfont}
\usepackage{xspace}
\usepackage{cancel}
\usepackage{mathtools}
\usepackage{amsthm}
\makeatletter
\@ifundefined{theorem}{\newtheorem{theorem}{Theorem}}{}
\@ifundefined{corollary}{\newtheorem{corollary}[theorem]{Corollary}}{}
\@ifundefined{lemma}{\newtheorem{lemma}[theorem]{Lemma}}{}
\@ifundefined{proposition}{\newtheorem{proposition}[theorem]{Proposition}}{}
\makeatother
\newcommand{\PiZeroE}{\Pi^{0,E}_k}
\title[Conforming/Non-conforming Virtual Elements and application t]{Conforming/Non-conforming Virtual Elements and application to elasticity problems in curved three-dimensional domains}
\author{L. Beir\~ao da Veiga, F. Dassi, A. Russo, M. Trezzi}
\date{Source: arXiv:2605.28041v1 (CC BY 4.0)}
\begin{document}
\maketitle

\noindent\emph{Source and licence.} Conforming/Non-conforming Virtual Elements and application to elasticity problems in curved three-dimensional domains. Version arXiv:2605.28041v1 (CC BY 4.0), distributed under the Creative Commons Attribution 4.0 International licence, as recorded in the arXiv licence metadata for this exact version and shown on the abstract page https://arxiv.org/abs/2605.28041v1. Full rights notice: licenses/arxiv-2605.28041-CC-BY-4.0.txt.\par\smallskip

\noindent\textit{Editorial scope.} Counted unit: the Proposition whose source label is \texttt{prop:judoboy} in the article (source0.tex lines 719--728, source label \texttt{prop:judoboy}), delivered together with the article's own complete original proof of it (source0.tex lines 729--775, one \texttt{proof} environment of 47 lines). No mathematics is written, repaired or re-derived: statement and proof are the article's own text line for line. Required context retained uncounted: sec:mesh-ass (lines 203--203); lem:w:prop (lines 601--612); cor:easy (lines 685--695); lm:beiraos-lemma (lines 697--704). Every definition, display and label of those lines is retained unchanged. Omitted from this package and not redistributed: 1--202, 204--600, 613--684, 696--696, 705--718, 776--1377. Nothing inside a retained excerpt is omitted or altered: every retained body is delivered exactly as the article's lines are written, so no omission receipt of any kind accompanies this package. Citations: the retained excerpts carry no citation command at all, so no bibliography entry of the article is transcribed and no reference is redistributed. Reference adaptation: none was needed and none was made; every reference inside the delivered unit resolves inside it, so no unresolved reference or citation is produced. Printed slips kept literal: no printed word, symbol, hypothesis or number of the counted statement or of its proof was altered. Layout and class: the article's own class is \texttt{article} with packages including geometry, amsmath, amsfonts, amsthm, amssymb, mathtools, cancel, comment, enumerate, graphicx, hyperref, lmodern, xcolor, authblk; the wrapper is the lane's pinned \texttt{amsart} editorial wrapper at 10pt on A4, which restates the theorem-like environments the retained text needs and adds the article's own macro definitions that the retained text needs (\texttt{\textbackslash PiZeroE}), with extra wrapper packages (bm, bbm, dsfont, xspace, cancel, mathtools). The delivered numbering is the unit's own. Assets: no retained span contains a figure environment, a graphics-inclusion command, a PDF-page inclusion, a table or a TikZ picture; no image file is copied into this package, the package declares no asset dependency and no image file is redistributed with it. Licence: version arXiv:2605.28041v1 of this article is distributed under the Creative Commons Attribution 4.0 International licence, as recorded in the arXiv licence metadata for this exact version and shown on the abstract page https://arxiv.org/abs/2605.28041v1. A full rights notice is delivered at licenses/arxiv-2605.28041-CC-BY-4.0.txt. Characters: every retained excerpt is pure ASCII, so the delivered engine, XeLaTeX with its default Latin Modern text font, reports no missing character.
\subsection{Basic notations and mesh assumptions}\label{sec:mesh-ass}

\begin{lemma}\label{lem:w:prop}
Let $v_h \in V^k_h(E)$ and $w \in H^1(E)$ be such that:
\begin{itemize}
    \item $w = v_h$ on $F$, for all $F \in \mathcal{F}^c_E$;
    \item $\Pi^{0,F}_{k-1}(w - v_h) = 0$ for all $F \in \mathcal{F}^{nc}_E$;
    \item $\PiZeroE(w -v_h) = 0$.
\end{itemize}
Then, the following estimate holds:
\[
| v_h |_{1,E} \leq | w |_{1,E} \, .
\]
\end{lemma}

\begin{corollary}\label{cor:easy}
Let the mesh assumptions in Section \ref{sec:mesh-ass} hold. 
Then, the following estimate holds for all $E \in \Omega_h$ and $v_h \in V_h^k(E)$:
\begin{equation}
\| \PiZeroE v_h \|_{0,E}
\lesssim
\| \Pi^{0,E}_{k-2} v_h \|_{0,E}
+ \sum_{F \in \mathcal{F}^{c}_E} h_E^{1/2} \| v_h \|_{0,F}
+ \sum_{F \in \mathcal{F}^{nc}_E} h_E^{1/2} \| \Pi^{0,F}_{k-1} v_h \|_{0,F} \, .
\end{equation}
\end{corollary}

\begin{lemma}\label{lm:beiraos-lemma}
Let the mesh assumptions in Section \ref{sec:mesh-ass} hold.
Let $v_h \in V_h^k(E)$, $E \in \Omega_h$. Then it exists a function $w \in H^1(E)$ that satisfies the first and second conditions in Lemma \ref{lem:w:prop}, plus the additional property
\begin{equation}\label{eq:equiv:br}
h_E |w|_{1,E} \simeq \| w \|_{0,E} \simeq \sum_{F \in \mathcal{F}^{c}_E} h_E^{1/2} \| v_h \|_{0,F}
+ \sum_{F \in \mathcal{F}^{nc}_E} h_E^{1/2} \| \Pi^{0,F}_{k-1} v_h \|_{0,F}  \, .
\end{equation}
\end{lemma}

\begin{proposition}\label{prop:judoboy}
Let the mesh assumptions in Section \ref{sec:mesh-ass} hold.
Then  the following estimate holds for all $E \in \Omega_h$ and $v_h \in V_h^k(E)$
\[
|v_h|_{1,E} 
\lesssim h_E^{-1} \| \Pi^{0,E}_{k-2} v_h \|_{0,E}
+ \sum_{F \in \mathcal{F}^{c}_E} h_E^{-1/2} \| v_h \|_{0,F}
+ \sum_{F \in \mathcal{F}^{nc}_E} h_E^{-1/2} \| \Pi^{0,F}_{k-1} v_h \|_{0,F}
\]
\end{proposition}

\begin{proof}
Given $B$, the ball contained in $E$ associated to the shape-regularity condition in Section \ref{sec:mesh-ass},
we denote by $\phi$ a bubble function with support in $B \subset E$ such that $\| \phi \|_{L^\infty(E)} = 1$.
Then, we select a polynomial $p \in \mathbb{P}_k(E)$ such that the function defined as $z = w + \phi p$ satisfies
\[
\int_E (v_h - z) q \, {\rm d }E = 0 \, , \qquad \forall q \in \mathbb{P}_k(E) \, ,
\]
where $w$ is the function introduced in Lemma \ref{lm:beiraos-lemma}.
%%
We note that, since the function $z$ satisfies the assumptions in Lemma \ref{lem:w:prop} by construction, it holds
\[
|v_h|_{1,E} \leq | z |_{1,E} \, .
\]
By standard polynomial properties and by definition of $z$, we have that
\[
\| p \|^2_{0,E}
\lesssim
\int_E \phi p^2 {\rm d}E
=
\int_E (\PiZeroE v_h - w) \, p \,  {\rm d}E \, ;
\]
therefore the Cauchy-Schwarz inequality gives
\[
\| p \|_{0,E} 
\leq 
\| \PiZeroE v_h\|_{0,E} + \| w \|_{0,E} \, .
\]
Exploiting Corollary \ref{cor:easy} and Lemma \ref{lm:beiraos-lemma}, we deduce
\[
\| p \|_{0,E} 
\lesssim
 \| \Pi^{0,E}_{k-2} v_h \|_{0,E}
+ \sum_{F \in \mathcal{F}^{c}_E} h_E^{1/2} \| v_h \|_{0,F}
+ \sum_{F \in \mathcal{F}^{nc}_E} h_E^{1/2} \| \Pi^{0,F}_{k-1} v_h \|_{0,F} \, .
\]
In conclusion, we obtain the desired estimate 
combining the above bounds with some manipulations and an inverse estimate for polynomials
\[
\begin{aligned}
|v_h|_{1,E} &\leq | z |_{1,E}
\leq | w |_{1,E} + | \phi p|_{1,E} \\
&\leq h_E^{-1} \| \Pi^{0,E}_{k-2} v_h \|_{0,E}
+ \sum_{F \in \mathcal{F}^{c}_E} h_E^{-1/2} \| v_h \|_{0,F}
+ \sum_{F \in \mathcal{F}^{nc}_E} h_E^{-1/2} \| \Pi^{0,F}_{k-1} v_h \|_{0,F} \, .
\end{aligned}
\]
 \end{proof}

\end{document}
